% ========== ASSEMBLING PROCESS ==========
% Building from primitives, component composition, and real-world examples
% Source: Symphony/Content/The Minimal IDE, The Assembling

\section*{The Assembling Process}
\addcontentsline{toc}{section}{The Assembling Process}

\lettrine{T}{he} true power of Symphony's minimal core and generic primitives becomes evident in the assembling process—how extensions combine these building blocks to create sophisticated functionality that rivals or exceeds traditional built-in features. This process demonstrates Symphony's philosophy of \textit{composition over configuration}.

\subsection*{Building from Primitives}
\addcontentsline{toc}{subsection}{Building from Primitives}

The assembling process follows a systematic approach where complex features emerge from the combination of simple, well-defined primitives:

\begin{table}[h]
\centering
\begin{tabular}{@{}p{0.25\textwidth}p{0.35\textwidth}p{0.35\textwidth}@{}}
\toprule
\textbf{Primitive Category} & \textbf{Available Primitives} & \textbf{Assembly Capabilities} \\
\midrule
\textbf{Communication} & 
JSON-RPC, HTTP, WebSocket, stdio, TCP & 
Any protocol implementation, custom message formats \\
\midrule
\textbf{Process Management} & 
Spawn, monitor, restart, resource control & 
Tool integration, service orchestration, pipeline management \\
\midrule
\textbf{UI Framework} & 
Components, panels, overlays, layouts & 
Custom interfaces, workflow-specific UIs, data visualization \\
\midrule
\textbf{Data Handling} & 
Serialization, validation, transformation & 
Protocol adaptation, data processing, format conversion \\
\midrule
\textbf{Event System} & 
Pub/sub, hooks, lifecycle events & 
Reactive programming, workflow automation, integration \\
\bottomrule
\end{tabular}
\caption{Primitive Categories and Assembly Capabilities}
\end{table}

\begin{infobox}[title=Assembly Philosophy]
The assembling process is guided by the principle that \textbf{simple primitives, when properly combined, can create arbitrarily complex functionality}. This approach ensures that extensions remain maintainable while achieving sophisticated capabilities.
\end{infobox}

\subsection*{Component Composition Patterns}
\addcontentsline{toc}{subsection}{Component Composition Patterns}

Symphony extensions follow established composition patterns that promote reusability and maintainability:

\subsubsection*{Layered Architecture Pattern}

Extensions typically organize functionality in layers:

\begin{compactlist}
    \item \textbf{Protocol Layer}: Handles communication with external tools
    \item \textbf{Service Layer}: Implements business logic and data processing
    \item \textbf{UI Layer}: Provides user interface and interaction
    \item \textbf{Integration Layer}: Connects with other extensions and core features
\end{compactlist}

\subsubsection*{Modular Composition Pattern}

Complex extensions can be composed of smaller, focused components that work together to provide comprehensive functionality while maintaining clear separation of concerns.

\subsection*{Real-World Assembly Examples}
\addcontentsline{toc}{subsection}{Real-World Assembly Examples}

The assembling process enables extensions to create sophisticated development tools:

\begin{description}[leftmargin=4cm,labelwidth=3.5cm]
    \item[\textbf{Language Support}] Combining language servers, syntax highlighting, and debugging tools
    \item[\textbf{Source Control}] Integrating Git operations, visual diff, and merge conflict resolution
    \item[\textbf{Testing Frameworks}] Orchestrating test runners, result visualization, and coverage reporting
    \item[\textbf{Build Systems}] Coordinating compilation, dependency management, and deployment
\end{description}

\subsection*{Assembly Benefits}
\addcontentsline{toc}{subsection}{Assembly Benefits}

The primitive-based assembly approach provides several key advantages:

\begin{successbox}
Extensions built through primitive assembly are more maintainable, testable, and adaptable than monolithic implementations. They can evolve independently while maintaining compatibility with the broader ecosystem.
\end{successbox}

\begin{compactlist}
    \item \textbf{Modularity}: Components can be developed, tested, and updated independently
    \item \textbf{Reusability}: Primitives can be combined in different ways for different use cases
    \item \textbf{Flexibility}: Extensions can adapt to new requirements without architectural changes
    \item \textbf{Community Innovation}: Developers can create novel combinations and patterns
    \item \textbf{Quality Assurance}: Smaller, focused components are easier to test and debug
\end{compactlist}

[composition]
# Define how components work together
language_server = "rust-analyzer-client"
build_system = "cargo-integration"
debugger = "rust-debugger"
test_runner = "rust-testing"

# Shared configuration
[shared_config]
rust_toolchain = "stable"
target_directory = "target"
\end{lstlisting}

\subsection*{Progressive Enhancement}
\addcontentsline{toc}{subsection}{Progressive Enhancement}

Symphony's assembling process supports progressive enhancement, where extensions can provide increasingly sophisticated functionality as more components become available:

\subsubsection*{Base Functionality}

Every extension starts with a minimal base that works with core primitives:

\begin{compactlist}
    \item \textbf{Basic Text Editing}: Leverages core editor primitives
    \item \textbf{File Operations}: Uses core file system integration
    \item \textbf{Simple UI}: Utilizes basic panel and component primitives
\end{compactlist}

\subsubsection*{Enhanced Functionality}

As additional extensions are installed, functionality can be enhanced:

\begin{compactlist}
    \item \textbf{Language Intelligence}: Integrates with language server extensions
    \item \textbf{Advanced UI}: Leverages sophisticated UI framework extensions
    \item \textbf{Tool Integration}: Connects with build system and debugger extensions
\end{compactlist}

\begin{successbox}
Progressive enhancement ensures that extensions provide value immediately while becoming more powerful as the user's extension ecosystem grows. This approach eliminates the "all-or-nothing" problem common in traditional IDEs.
\end{successbox}

\subsection*{Real-World Example: LSP Support}
\addcontentsline{toc}{subsection}{Real-World Example: LSP Support}

To illustrate the assembling process, let's examine how Language Server Protocol (LSP) support is built from Symphony's primitives—a feature that would be hardcoded in traditional IDEs.

\subsubsection*{LSP Extension Architecture}

The LSP extension demonstrates sophisticated assembly from basic primitives:

\begin{lstlisting}[language=javascript, caption=LSP Extension Assembly Process]
class LanguageServerExtension {
    constructor() {
        // Assemble communication primitives
        this.processManager = new ProcessManager();
        this.jsonRpcClient = new JsonRpcClient();
        this.messageQueue = new MessageQueue();
        
        // Assemble UI primitives
        this.diagnosticsPanel = new DiagnosticsPanel();
        this.completionProvider = new CompletionProvider();
        this.hoverProvider = new HoverProvider();
        
        // Assemble integration primitives
        this.editorIntegration = new EditorIntegration();
        this.fileWatcher = new FileWatcher();
    }
    
    async activate() {
        // 1. Process Management Assembly
        this.languageServer = await this.processManager.spawn({
            command: this.getLanguageServerCommand(),
            communication: 'stdio',
            protocol: 'json-rpc'
        });
        
        // 2. Communication Assembly
        this.client = this.jsonRpcClient.connect(this.languageServer);
        await this.initializeLanguageServer();
        
        // 3. UI Assembly
        this.setupDiagnosticsUI();
        this.setupCompletionUI();
        this.setupHoverUI();
        
        // 4. Integration Assembly
        this.setupEditorIntegration();
        this.setupFileWatching();
    }
    
    async initializeLanguageServer() {
        // Use communication primitives for LSP initialization
        const initResult = await this.client.request('initialize', {
            processId: process.pid,
            rootUri: symphony.workspace.rootUri,
            capabilities: this.getClientCapabilities()
        });
        
        this.serverCapabilities = initResult.capabilities;
        await this.client.notify('initialized', {});
    }
    
    setupDiagnosticsUI() {
        // Assemble UI primitives for diagnostics display
        this.client.onNotification('textDocument/publishDiagnostics', 
            (params) => {
                this.diagnosticsPanel.update(params.uri, params.diagnostics);
                this.editorIntegration.showDiagnostics(params);
            }
        );
    }
    
    setupCompletionUI() {
        // Assemble completion provider from primitives
        this.editorIntegration.registerCompletionProvider({
            triggerCharacters: ['.', ':', '>'],
            provideCompletions: async (document, position) => {
                const result = await this.client.request(
                    'textDocument/completion',
                    {
                        textDocument: { uri: document.uri },
                        position: position
                    }
                );
                return this.completionProvider.format(result);
            }
        });
    }
}
\end{lstlisting}

\subsubsection*{Primitive Composition Breakdown}

The LSP extension demonstrates how multiple primitive categories work together:

\begin{expandedlist}
    \item \textbf{Process Primitives}: Spawn and manage the language server process
    \item \textbf{Communication Primitives}: Handle JSON-RPC protocol communication
    \item \textbf{UI Primitives}: Create diagnostics panels, completion popups, hover tooltips
    \item \textbf{Editor Primitives}: Integrate with text editing, cursor management, selection
    \item \textbf{File System Primitives}: Watch for file changes, handle document synchronization
    \item \textbf{Event Primitives}: Coordinate between different components and handle lifecycle
\end{expandedlist}

\subsubsection*{Configuration Assembly}

The extension's manifest shows how primitives are declared and assembled:

\begin{lstlisting}[language=toml, caption=LSP Extension Manifest]
[extension]
name = "language-server-client"
version = "2.1.0"
description = "Generic Language Server Protocol client"

# Declare required primitives
[capabilities]
process_management = true
editor_integration = true
ui_modification = true
file_system_access = true

# Communication primitive configuration
[protocols.lsp]
name = "language-server-protocol"
transport = ["stdio", "tcp"]
message_format = "json-rpc"
bidirectional = true

# UI primitive configuration
[ui_components]
[[ui_components.panels]]
id = "diagnostics"
title = "Problems"
location = "bottom"
icon = "alert-circle"

[[ui_components.providers]]
type = "completion"
languages = ["*"]
trigger_characters = ["."]

# Process primitive configuration
[processes.language_server]
command = "${config.languageServer.command}"
args = "${config.languageServer.args}"
working_directory = "${workspaceRoot}"
auto_restart = true
health_check_interval = 30

# Integration primitive configuration
[editor_integration]
document_sync = true
diagnostics = true
completion = true
hover = true
signature_help = true
\end{lstlisting}

\subsection*{Assembly Benefits and Outcomes}
\addcontentsline{toc}{subsection}{Assembly Benefits and Outcomes}

The assembling process provides several key benefits over traditional hardcoded approaches:

\begin{table}[h]
\centering
\begin{tabular}{@{}p{0.45\textwidth}p{0.45\textwidth}@{}}
\toprule
\textbf{Traditional Hardcoded LSP} & \textbf{Symphony Assembled LSP} \\
\midrule
Fixed to specific LSP version & Adapts to any LSP version automatically \\
Limited customization options & Fully customizable UI and behavior \\
Vendor-controlled feature set & Community-driven feature evolution \\
Monolithic implementation & Modular, reusable components \\
Difficult to extend or modify & Easy to extend and customize \\
Single language server per language & Multiple servers, user choice \\
\bottomrule
\end{tabular}
\caption{Comparison: Hardcoded vs Assembled LSP Implementation}
\end{table}

\begin{alertbox}
The LSP example demonstrates that Symphony's assembled extensions can not only match the functionality of traditional built-in features but often exceed them in flexibility, customization, and innovation potential.
\end{alertbox}

\subsubsection*{Real-World Impact}

This assembling approach has enabled the Symphony community to create:

\begin{compactlist}
    \item \textbf{Multi-Language Support}: Single extension supporting multiple language servers
    \item \textbf{Custom Protocols}: Extensions for proprietary or experimental language servers
    \item \textbf{Enhanced UIs}: Sophisticated debugging interfaces, advanced completion systems
    \item \textbf{Workflow Integration}: LSP integration with build systems, testing frameworks, and deployment tools
\end{compactlist}

The assembling process transforms Symphony from a minimal IDE into a powerful, infinitely extensible development platform where the community's creativity is the only limit to what can be achieved.